% !TEX root = ./main.tex
% \documentclass[11pt, oneside]{article} 
\usepackage[T1]{fontenc}
\usepackage[ngerman]{babel}	
\usepackage[default]{sourcesanspro}
\usepackage{eulervm}
	\renewcommand{\familydefault}{\sfdefault}
\usepackage[a4paper, margin=2.54cm]{geometry}
\usepackage{changepage}
    \newenvironment{hanging}{\begin{adjustwidth}{1.5\baselineskip}{}}{\end{adjustwidth}}
\usepackage{fancyhdr}
\usepackage[dvipsnames]{xcolor}
\usepackage{ragged2e}
	\justifying
\usepackage{mathtools, amssymb}
\usepackage{stmaryrd}
\usepackage[shortlabels, inline]{enumitem}
	\setlist[enumerate,1]{label={(\alph*)}}
\usepackage{tikz,ifthen}
	\usetikzlibrary{cd, backgrounds, babel, 
	calc, intersections, shadows}
\usepackage{amsthm,thmtools, etoolbox}
	\theoremstyle{definition} 
		\newtheorem{theorem}{Satz}
		\numberwithin{theorem}{section}
		\newtheorem{proposition}[theorem]{Proposition}
		\newtheorem{lemma}[theorem]{Lemma}
		\newtheorem{corollary}[theorem] {Korollar}
		\newtheorem{corollary_weak}[theorem] {Korollar}
		\newtheorem{definition}[theorem] {Definition}
		\newtheorem{notation}[theorem]{Notation}
		\newtheorem{convention}[theorem]{Konvention}
		\newtheorem{example}[theorem]{Beispiel}
		\newtheorem{construction}[theorem]{Konstruktion}
	\theoremstyle{remark}
		\newtheorem{remark} [theorem] {Bemerkung}
    \newcounter{claimcounter}
        \newtheorem{claim}[claimcounter]{Aussage}
        \renewcommand{\theclaim}{\arabic{claimcounter}}
        \BeforeBeginEnvironment{proof}{%
         \setcounter{claimcounter}{0}%
        }
        % \newenvironment{claim}[1]{\par\noindent\textit{Aussage:}\space#1}{}
        \newenvironment{proofofclaim}[1]{
           % \begin{hanging}
                \par\textit{\small Beweis. }\space#1}{\leavevmode\unskip\penalty9999 \hbox{}\nobreak\hfill\quad\hbox{\color{OliveGreen}{$\blacksquare$}}
           % \end{hanging} 
		   	\linebreak
        }
		

		\renewenvironment{proof}{\textit{Beweis.}}{\leavevmode\unskip\penalty9999 \hbox{}\nobreak\hfill\quad\hbox{$\square$}}

\usepackage[most]{tcolorbox}
	%footnote at end of page 
	\newif\iffootaux
			\newcounter{footaux}% for \csname
			\newcounter{auxpage}% to handle frontmatter etc.
			\setcounter{auxpage}{1}

			\makeatletter
			\tcbset{every box/.style={before upper app=\tcb@footnote@setup}}

			\def\tcb@footnote@setup{\let\footnotetext=\footaux
			\let\footnote=\footnoteaux}
			
			\def\footnoteaux{\@ifnextchar[\@xfootnoteaux{\stepcounter{footnote}%
			\footnotemark[\thefootnote]\@footaux[\thefootnote]}}
			
			\def\@xfootnoteaux[#1]{\footnotemark[#1]\@footaux[#1]}

			\def\footaux{\@ifnextchar[\@footaux{\@footaux[\thefootnote]}}
			\def\@footaux[#1]#2{\protected@write{\@auxout}{\let\theauxpage\relax}{\string\newfootnote{\theauxpage}{#1}{#2}}}

			\newcommand{\newfootnote}[3]{% #1 = \theauxpage, #2 = \thefoontnote, #3 = text
			\ifnum#1=1\relax
				\footnotetext[#2]{#3}%
			\else
				\stepcounter{footaux}%
				\expandafter\xdef\csname footauxpage\thefootaux\endcsname{#1}%
				\expandafter\xdef\csname footauxnumber\thefootaux\endcsname{#2}%
				\expandafter\protected@xdef\csname footauxtext\thefootaux\endcsname{#3}%
			\fi}

			% to run at the end of every page (debug version)
			\newcommand{\addfootaux}{\loop
			\@ifundefined{footauxpage\thefootaux}{\footauxfalse}{%
				\count1=\csname footauxpage\thefootaux\endcsname\relax
				\advance\count1 by -1
				\ifnum\count1=\value{auxpage}%
				\expandafter\let\expandafter\footauxnumber\csname footauxnumber\thefootaux\endcsname
				\expandafter\let\expandafter\footauxtext\csname footauxtext\thefootaux\endcsname
				\footnotetext[\footauxnumber]{\footauxtext}%
				\footauxtrue
				\else
				\footauxfalse
				\fi}
			\iffootaux
				\global\expandafter\let\csname footauxpage\thefootaux\endcsname\relax
				\global\expandafter\let\csname footauxnumber\thefootaux\endcsname\relax
				\global\expandafter\let\csname footauxtext\thefootaux\endcsname\relax
				\stepcounter{footaux}%
			\repeat
			\stepcounter{auxpage}}
			\makeatother

			\AtBeginDocument{\setcounter{footaux}{1}}% reset counter after \newfootnote
		\AddToHook{shipout/after}{\addfootaux}


	\newtcbox{\inlinebox}[1][]{enhanced,
		box align=base,
		nobeforeafter,
		colback=white,
		frame hidden,
		size=small,
		left=0pt,
		right=0pt,
		boxsep=2pt,
		rounded corners,
		#1
	}
	\renewcommand{\thempfootnote}{\fnsymbol{mpfootnote}}
	\tcbuselibrary{breakable}
		\newtcolorbox{idea}{box align = top, breakable, enhanced jigsaw,
		colback=green!5!white,
		colframe=green!75!black,fonttitle=\bfseries,
		title={Idee}
		}
		\newtcolorbox{Bastelanleitung}{box align = top, breakable, enhanced jigsaw,
		colback=green!5!white,
		colframe=green!75!black,fonttitle=\bfseries,
		title={Bastelanleitung}
		}
		\newtcolorbox{marked}{box align = top,
		boxrule=0pt,
		frame hidden,
		sharp corners,
		enhanced,
		borderline west={3pt}{0pt}{OliveGreen},
		colback = white,
		fonttitle=\bfseries,
		}	
	\tcbuselibrary{theorems}
		\newcommand{\mytcolorboxenvironment}[2]{
			\tcolorboxenvironment{#1}{
			box align = top, 
			breakable, 
			enhanced jigsaw,
			colback= #2!5!white,
			colframe=#2,
			before skip=10pt,after skip=10pt,
			}
		}
		\newcommand{\mytcolorboxenvironmentMarked}[2]{
			\tcolorboxenvironment{#1}{
				breakable, enhanced jigsaw,
				boxrule=0pt,
				frame hidden,
				sharp corners,
				enhanced,
				borderline west={3pt}{0pt}{#2},
				colback = white,
				fonttitle=\bfseries,
			}
		}
		\newcommand{\mytcolorboxenvironmentNoCorners}[2]{
			\tcolorboxenvironment{#1}{
				breakable, enhanced jigsaw,
				boxrule=0pt,
				frame hidden,
				sharp corners,
				enhanced,
				colback = #2,
				fonttitle=\bfseries,
			}
		}
		\mytcolorboxenvironment{definition}{cyan}
		\mytcolorboxenvironment{notation}{blue}
		\mytcolorboxenvironment{convention}{blue}
		\mytcolorboxenvironment{lemma}{Lavender}
		\mytcolorboxenvironment{proposition}{orange}
		\mytcolorboxenvironment{theorem}{red}
		\mytcolorboxenvironment{corollary}{RedOrange}
		\mytcolorboxenvironment{corollary_weak}{Lavender}
		\mytcolorboxenvironmentMarked{remark}{OliveGreen}
		\mytcolorboxenvironmentNoCorners{proof}{gray!5!white}
		\mytcolorboxenvironmentMarked{claim}{OliveGreen!100!white}
\usepackage[colorlinks, allcolors=blue]{hyperref}
\usepackage[german, nameinlink, noabbrev]{cleveref}
	\newcommand{\myref}[1]{\cref{#1} (\textcolor{gray}{\footnotesize {\nameref*{#1}}})}
    \crefname{proposition}{Proposition}{Propositionen}
    \crefname{claim}{Aussage}{Aussageen}
\usepackage{csquotes}
\usepackage[backend=biber, style=alphabetic, maxcitenames=1, natbib]{biblatex}
	%\usepackage{filecontents}
	\begin{filecontents}[overwrite]{Bachelorarbeit.bib}
		@BOOK{ahlfors,
			AUTHOR = {Ahlfors, L. V.},
			PUBLISHER = {McGraw-Hill Book Company},
			DATE = {1966},
			EDITION = {2},
			SUBTITLE = {An introduction to the theory of analytic functions of one complex variable},
			TITLE = {Complex Analysis},
		}
		@book{hatcher,
			address = {New York},
			title = {Algebraic topology},
			%isbn = {978-0-521-79540-1},
			%language = {eng},
			publisher = {Cambridge university press},
			author = {Hatcher, Allen},
			year = {2001},
		}
		@book{forster,
  			title={Lectures on Riemann Surfaces},
  			author={Forster, O. and Gilligan, B.},
			%isbn={9781461259619},
			%series={Graduate Texts in Mathematics},
			%url={https://books.google.de/books?id=6wvpBwAAQBAJ},
			year={2012},
			publisher={Springer New York}
		}
		@book{lee2013introduction,
			title={Introduction to Smooth Manifolds},
			author={Lee, J.M.},
			%isbn={9781475756012},
			series={Graduate Texts in Mathematics},
			%url={https://books.google.de/books?id=A0sSswEACAAJ},
			year={2013},
			publisher={Springer New York}
		}
	\end{filecontents}
	\addbibresource{Bachelorarbeit.bib}
\usepackage{wrapfig}
\usepackage{scalerel}
\usepackage{stackengine,wasysym}
	%reallywidetilde
	\newcommand\reallywidetilde[1]{\ThisStyle{%
	\setbox0=\hbox{$\SavedStyle#1$}%
	\stackengine{-.1\LMpt}{$\SavedStyle#1$}{%
		\stretchto{\scaleto{\SavedStyle\mkern.2mu\AC}{.5150\wd0}}{.6\ht0}%
	}{O}{c}{F}{T}{S}%
	}}


\newcommand*{\nsection}[1]{
    \section*{#1}
    \addcontentsline{toc}{section}{#1}
}
\newcommand*{\nsubsection}[1]{
    \subsection*{#1}
    \addcontentsline{toc}{subsection}{#1}
}
\renewcommand{\subset}{\subseteq}
\newcommand{\N}{\mathbb{N}}
\newcommand{\Z}{\mathbb{Z}}
\newcommand{\Q}{\mathbb{Q}}
\newcommand{\R}{\mathbb{R}}
\newcommand{\C}{\mathbb{C}}
\newcommand*{\simsym}{\mathord\sim}
\newcommand{\mymarginpar}[1]{\marginpar{\tiny #1}}
\DeclareMathOperator{\Hom}{Hom}
\DeclareMathOperator{\End}{End}
\DeclareMathOperator{\id}{id}
\DeclareMathOperator{\ev}{ev}
\DeclareMathOperator{\Image}{Im}
\NewCommandCopy{\ordinaryforall}{\forall}
\NewCommandCopy{\ordinaryexists}{\exists}
\RenewDocumentCommand{\forall}{}{\mathop{{}\ordinaryforall}}
\RenewDocumentCommand{\exists}{}{\mathop{{}\ordinaryexists}}
\NewDocumentCommand{\existsunique}{}{\mathop{\ordinaryexists!}}
\newcommand{\norm}[1]{\left\lvert#1\right\rvert}
\newcommand{\Norm}[1]{\left\lVert#1\right\rVert}
\newcommand{
	\explainabove
}[2]{
	\overset{
	 	\mathclap{
			\tikz 
				\node [gray] {$\downarrow$} 
					node [above=0.5ex, gray, font=\tiny] {#1};
		}
	}{#2}
}
\newcommand{
	\explainbelow
}[2]{
       	 \underset{
	 	\mathclap{
			\tikz 
				\node [gray] {$\uparrow$} 
					node [below=0.5ex, gray, font=\tiny] {#1};
		}
	}{#2}
}

\newcommand{
	\explainabovebelow
}[3]{
	\explainabove{#1}{\explainbelow{#2}{#3}}
}
\newcommand{\arrow}{\rightarrow}
\newcommand{\longarrow}{\longrightarrow}
\setlength{\parindent}{0pt}
